%% =================================================================
%% Kingma, Ba. "Adam: A Method for Stochastic Optimization."
%% ICLR 2015. arXiv:1412.6980v9.
%%
%% Reconstruction of page 12 (printed p. 12) as study notes.
%% Generated by: reproduce-all-pages-basic-tex (batch 1-15)
%%
%% Structure: Sec 10 APPENDIX + Sec 10.1 CONVERGENCE PROOF with
%% Definition 10.1, Lemma 10.2, Lemma 10.3 (proved by induction).
%% All display equations are UNNUMBERED — no equation counter bump.
%% Anomaly preserved: the source inconsistently writes g_T and g_{T,i}
%% in the same derivation (e.g., "||g_{1:T,i}||^2 - g_T^2" in one line
%% and "||g_{1:T,i}||^2 - g_{T,i}^2" in the next). Reproduced verbatim.
%% =================================================================

\documentclass[11pt]{article}

\usepackage[margin=1in]{geometry}
\usepackage{amsmath, amssymb}
\usepackage{mathrsfs}
\usepackage{bm}
\usepackage{xcolor}
\usepackage{fancyhdr}

\pagestyle{fancy}
\fancyhf{}
\lhead{\small Published as a conference paper at ICLR 2015}
\rhead{}
\cfoot{\small 12 $\to$ \arabic{page}}
\renewcommand{\headrulewidth}{0.4pt}

\setcounter{page}{1}
\setcounter{equation}{12}

\begin{document}

\section*{10 \textsc{Appendix}}

\subsection*{10.1 \textsc{Convergence proof}}

\noindent\textbf{Definition 10.1.} \emph{A function
$f : \mathbb{R}^d \to \mathbb{R}$ is convex if for all
$x, y \in \mathbb{R}^d$, for all $\lambda \in [0, 1]$,}
\[
\lambda f(x) + (1 - \lambda) f(y) \geq f(\lambda x + (1 - \lambda) y)
\]

\noindent
Also, notice that a convex function can be lower bounded by a
hyperplane at its tangent.

\noindent\textbf{Lemma 10.2.} \emph{If a function
$f : \mathbb{R}^d \to \mathbb{R}$ is convex, then for all
$x, y \in \mathbb{R}^d$,}
\[
f(y) \geq f(x) + \nabla f(x)^\mathsf{T}(y - x)
\]

\noindent
The above lemma can be used to upper bound the regret and our
proof for the main theorem is constructed by substituting the
hyperplane with the Adam update rules.

The following two lemmas are used to support our main theorem.
We also use some definitions simplify our notation, where
$g_t \triangleq \nabla f_t(\theta_t)$ and $g_{t,i}$ as the
$i^\text{th}$ element. We define $g_{1:t,i} \in \mathbb{R}^t$ as
a vector that contains the $i^\text{th}$ dimension of the
gradients over all iterations till $t$,
$g_{1:t,i} = [g_{1,i}, g_{2,i}, \ldots, g_{t,i}]$.

\noindent\textbf{Lemma 10.3.} \emph{Let $g_t = \nabla f_t(\theta_t)$
and $g_{1:t}$ be defined as above and bounded,
$\|g_t\|_2 \leq G$, $\|g_t\|_\infty \leq G_\infty$. Then,}
\[
\sum_{t=1}^{T} \sqrt{\frac{g_{t,i}^2}{t}} \leq 2 G_\infty \|g_{1:T,i}\|_2
\]

\noindent\emph{Proof.} We will prove the inequality using
induction over $T$.

The base case for $T = 1$, we have
$\sqrt{g_{1,i}^2} \leq 2 G_\infty \|g_{1,i}\|_2$.

For the inductive step,
\begin{align*}
\sum_{t=1}^{T} \sqrt{\frac{g_{t,i}^2}{t}}
 &= \sum_{t=1}^{T-1} \sqrt{\frac{g_{t,i}^2}{t}} + \sqrt{\frac{g_{T,i}^2}{T}} \\
 &\leq 2 G_\infty \|g_{1:T-1,i}\|_2 + \sqrt{\frac{g_{T,i}^2}{T}} \\
 &= 2 G_\infty \sqrt{\|g_{1:T,i}\|_2^2 - g_T^2} + \sqrt{\frac{g_{T,i}^2}{T}}
\end{align*}

\noindent
From,
$\|g_{1:T,i}\|_2^2 - g_{T,i}^2 + \dfrac{g_{T,i}^4}{4\|g_{1:T,i}\|_2^2}
 \geq \|g_{1:T,i}\|_2^2 - g_{T,i}^2$,
we can take square root of both side and have,
\begin{align*}
\sqrt{\|g_{1:T,i}\|_2^2 - g_{T,i}^2}
 &\leq \|g_{1:T,i}\|_2 - \frac{g_{T,i}^2}{2 \|g_{1:T,i}\|_2} \\
 &\leq \|g_{1:T,i}\|_2 - \frac{g_{T,i}^2}{2 \sqrt{T G_\infty^2}}
\end{align*}

\noindent
Rearrange the inequality and substitute the
$\sqrt{\|g_{1:T,i}\|_2^2 - g_{T,i}^2}$ term,
\[
G_\infty \sqrt{\|g_{1:T,i}\|_2^2 - g_T^2}
 + \sqrt{\frac{g_{T,i}^2}{T}}
 \leq 2 G_\infty \|g_{1:T,i}\|_2
\]
\hfill$\square$

% [page ends here — continues on page 13]

\end{document}
